\documentclass[instructions.tex]{subfiles}
\begin{document}
 
\chapter{On doit respecter les prêtres et les personnes consacrées à Dieu.}
 
\primo Si l'on doit du respect aux officiers d'un roi mortel, à plus forte raison doit-on respecter les ministres de Dieu. Celui qui vous écoute m'écoute, disoit Jésus-Christ à ses disciples, et celui qui vous méprise me méprise\footnote{Qui vos audit , me audit; qui vos spernit, me spernit. Luc. 10. 16.}. Manquer de respect aux prêtres, c'est en manquer à Dieu. Ne regardez pas leur personne, mais leur caractère, qui est redoutable aux anges mêmes, dit saint Grégoire.
 
Caractère si auguste, dit saint Chrysostome, qu'il est au-dessus de la pourpre et de la dignité royale, parce qu'il donne un pouvoir que les rois et les anges même n'ont pas. Remettre les péchés, lier ou délier les consciences, produire le corps et le sang de Jésus-Christ, l'offrir en sacrifice, le distribuer aux fidèles, donner la grâce par les sacremens aux vivans et aux mourans, chasser les démons, voilà ce que ne peuvent tous les rois, et voilà ce que peuvent tous les prêtres du Seigneur.
 
Par le sacré sacerdoce, ils sont d'ailleurs par état les médiateurs entre Dieu et les hommes, destinés à annoncer l'Évangile et la loi de Dieu à toute créature, aux puissances mêmes de la terre; ils sont les lieutenans de Dieu; mépriser leur sacré ministère, c'est s'en prendre à lui et le toucher à la prunelle de l'œil. Prenez garde, dit le Seigneur, de toucher à mes oints, à ceux qui me sont consacrés\footnote{Nolite tangere christos meos, et in prophetis meis nolite malignari. Ps. 104. 15.}. C'est pourquoi le Saint-Esprit nous avertit de baisser la tête devant les grands du monde et d'humilier notre âme devant un prêtre\footnote{Presbytero humilia animam tuam, et magnato humilia caput tuum 1. Eccli. 4. 7.}.
 
\secundo Mépriser les prêtres, les religieux et les pasteurs, est ordinairement une marque qu'on n'aime ni Dieu, ni la religion, ni son devoir. Les pharisiens se moquoient de Jésus-Christ, le souverain prêtre, parce qu'ils étoient avares, et aujourd'hui on se moque des prêtres, parce qu'on est vicieux. Un pasteur qui fait son devoir, a presque autant d'ennemis qu'il a de paroissiens qui vivent mal; aussitôt qu'il veut s'opposer aux désordres, empêcher les débauches, les intrigues et les injustices, on ne peut plus le souffrir.
 
Lorsqu'un pasteur vous reprend, il fait son devoir, parce qu'il est obligé de veiller sur votre conduite, et doit en répondre à Dieu. S'il vous reprend publiquement, prenez-vous-en à vous-même; pourquoi vos désordres sont-ils publics? Le berger doit-il se taire quand il voit le loup ravager son troupeau? S'il vous avertit des désordres de vos enfans, de vos domestiques ou de vos sujets, c'est un devoir que Dieu lui impose, et un service qu'il vous rend; quand vous le trouverez mauvais, c'est une marque que vous êtes ennemi de la vérité.
 
Le démon met tout en œuvre pour prévenir les fidèles contre les pasteurs; gardez-vous de donner dans ce piège; considérez votre pasteur comme un guide que Dieu vous a destiné pour vous conduire dans le chemin du ciel. S'il n'y marchoit pas lui-même, ce seroit un grand malheur à lui; mais, quand il ne seroit pas irréprochable, il n'en est pas moins votre pasteur: regardez son autorité, et non pas ses défauts; ayez confiance en lui pendant la vie, vous en aurez besoin à la mort.
 
\tertio Refuser aux pasteurs de l'Église le temporel qui leur est dû, est une injustice et une criante ingratitude. Saint Paul ne dit-il pas que celui qui sert à l'autel, doit vivre de l'autel, et qu'après avoir donné le spirituel aux autres, il est juste qu'il en reçoive au moins le temporel?
 
Comprend-on combien la charge d'un pasteur est pénible; les sollicitudes qui le dévorent, le compte qu'il rendra à Dieu des âmes qui lui sont confiées; les peines d'ouïr les confessions, d'administrer les sacremens, d'annoncer la divine parole, d'assister les pauvres, de visiter les malades, de les secourir en tout temps, la nuit, le jour, même au péril de sa santé et de sa vie; les soins qu'il doit prendre, sous peine d'en répondre, pour régler les consciences, pour empêcher le libertinage et les scandales, pour maintenir la piété, la paix dans les familles, et la crainte de Dieu parmi la jeunesse, etc.?
 
N'est-il pas juste que tant de travaux soient récompensés? Mais, loin d'user de reconnoissance, souvent on lui dispute ses droits, ou, comme Caïn, on lui donne ce qu'on a de plus mauvais. On paie libéralement l'honoraire à un jurisconsulte, à un avocat, tandis qu'on ne paie qu'à regret ce qu'on doit à un pasteur pour des peines bien plus grandes. Ce que vous refusez aux ministres de Dieu, dit saint Jean Chrysostome, ne vous profitera point, et vous sera enlevé par un injuste usurpateur\footnote{Quod noluisti dare sacerdoti, dabis impio militi. S. Chrys.}. Un pasteur ne doit pas attendre sa récompense des hommes, mais de Dieu seul. Jésus-Christ, le prince des pasteurs, n'a reçu des hommes sur la terre, pour reconnoissance de ses travaux, que des ingratitudes.
 
\section{Résolutions}
 
\primo Je regarderai les prêtres comme les lieutenans de Dieu, les dépositaires de sa puissance pour son culte et le salut des âmes.
 
\secundo Je les écouterai comme ses envoyés; quand ils m'instruiront, me donneront des avertissemens, et me reprendront de mes défauts et de mes vices.
 
\tertio Et jamais je ne m'opposerai imprudemment ni par malice, au bien qu'ils voudront faire dans leurs paroisses.

\prayer{Mon Dieu, donnez-nous des pasteurs selon votre cœur, et répandez vos grandes bénédictions sur leurs pénibles travaux.}
 
\end{document}
